\begin{lemma}[A globally maximum tree is fixed-pair maximal]
Let $(T,e_0,\varphi)$ be a globally maximum $k$-triple.  Then $T$ is
maximal among the Tashkinov trees for the fixed edge $e_0$ and fixed
coloring $\varphi$: if such a tree $S$ contains $V(T)$, then
$V(S)=V(T)$ (and in particular $V(S)\subseteq V(T)$).
\end{lemma}
\begin{proof}
By \noderef{GlobalMaximumTriple}, $T$ is a Tashkinov tree and
$|V(T)|=t(G)$.  The same node also supplies the $k$-triple condition, so
\noderef{KTriple} gives $\chi'(G)=k+1$ and hence
$k=\chi'(G)-1$.

Let $S$ be a Tashkinov tree for the same $e_0$ and $\varphi$, and suppose
$V(T)\subseteq V(S)$.  With $k=\chi'(G)-1$, the data
$(e_0,\varphi,S)$ satisfy exactly the predicate defining the candidate
orders in \noderef{TashkinovOrder}.  Thus $|V(S)|\leq t(G)=|V(T)|$.
On the other hand, the assumed inclusion gives $|V(T)|\leq |V(S)|$.
The two finite sets therefore have equal cardinality, and the inclusion
is an equality.  Hence $V(S)\subseteq V(T)$.  Together with the
Tashkinov-tree clause for $T$, this is precisely every clause of
\noderef{FixedPairMaximalTashkinovTree}.
\end{proof}
